\section{Weighted minimizing cuts and the forms on retained sides}\label[appendix]{sec:cut}

Prescribed-curvature perimeter minimization and convergence of its
boundaries have classical precedents in Massari and Massari--Pepe
\citep{massari1975,massariPepe1975}. For scalar-curvature applications,
prescribed-mean-curvature minimization is the $\mu$-bubble method
described by \citet[\S\(5\frac56\)]{gromov1996}.
We prove here the weighted,
varying-metric and product-end versions required by the cutting
construction, including the forms on both retained sides.

Fix a finite stage $r$ and the label $j_0$ to be processed.  The
packet and profile hypotheses are those of \cref{packet:category,packet:invariant}.
For the finite-horizon formulation in this appendix, only finitely many
labels and finitely many nonunit old warps occur in each incident diagram
under consideration.  The processed potential has a positive lower bound
\[
                         \Pi_{j_0}\ge\pi_{\rm cut}>0.
\]
A common lower bound for all labels at all stages is not required.
The constants satisfy $a_r\le .63$, $b_r\ge1.07$, and
$\lambda_r>\lambda_{r+1}$.  The extension to a countable collection of
labels is proved in the main text; every argument below uses only the
current pressure and the indicated finite lists.

Throughout, $E$ denotes a parent packet of base dimension $m\ge2$.
After the radial split it has the compact truncations, finite incidence,
and exact product ends of \cref{packet:category}.  At the initial split
we instead work on the complete original base, and use the locally finite
compact radial barriers of \cref{packet:initialization}.  We assume no
volume growth or bounded geometry on that initial base.

On a finite incident diagram, insert unit warps where necessary so that
the old warp lists agree.  Their metric is
\[
 \widehat g=g+\sum_{i=1}^q f_i^2\,\dd\theta_i^2,
 \qquad \omega=\prod_{i=1}^q f_i,
 \qquad \dd\mu=\omega\,\dd\vol_g.
\]
Empty products are one.  Fiber volume constants are suppressed throughout.
Adding a unit warp changes neither $\omega$ nor any formula below.
At the initial radial stage the list need only be finite on each compact
window; the compatible local expressions then agree on overlaps.  No
infinite-dimensional torus is used.  In particular all minimizing sets,
BV compactness, tangent cones, and singular-dimension bounds concern the
$m$-dimensional base, regardless of the length of the warp list.
Write $R=\Scal_{\widehat g}$ and
$\Div_\omega Y=\omega^{-1}\Div_g(\omega Y)$.
For a set $A$ of locally finite perimeter define its weighted perimeter by
\[
 P_\omega(A;U)=\int_{U\cap\partial^*A}\omega\,\dd\mathcal H_g^{m-1}.
\]
The reduced-boundary representation is the framework developed by
De Giorgi \citep{deGiorgi1954,deGiorgi1955}; the current formulation
and slicing originate in Federer and Fleming \citep{federerFleming1960}.
We use the precise representation and slicing statements in
\citet[Chapter~3, Section~4, and Chapter~6, Section~4]{simonGMT}.
All set equalities and inclusions involving measurable phases are understood
up to $\mu$-null sets.  The phase representative used for geometry will be
specified below.  Constants in local estimates may depend on a compact
window; after the radial split they can be chosen uniformly on each fixed
packet, without any uniformity between packets.  No scalar-curvature lower bound
is used to construct the cuts or to prove their local regularity.  The
parent scalar bound enters only at the end of this appendix, and only on
the required region $D_E$ of \cref{packet:required-set}.

\subsection{Smooth finite pressures and absolute minimizers}

\begin{lemma}[Realizing the processed profile]\label{cut:pressure}
There is a smooth finite compatible function $p_0$ and a number $c_*>0$
such that every $p=p_0+c$, $|c|<c_*$, satisfies
\begin{equation}\label{cut:pressure-gradient}
             2|\grad p|\le .67p^2+1.02\Pi_{j_0}.
\end{equation}
There are compatible smooth barrier domains $D^+,D^-$ whose interiors
contain the corresponding forced-core closures and whose closures lie in
the respective pole interiors.  Each admits a smooth vector field $Y^\pm$
with $|Y^\pm|\le1$, equal to the outward unit normal on its boundary, such
that, for all these shifts,
\begin{equation}\label{cut:barrier-pressure}
 p\ge1+\Div_\omega Y^+\quad\hbox{on }D^+,
 \qquad
 -p\ge1+\Div_\omega Y^-\quad\hbox{on }D^-.
\end{equation}
After the radial split all this data has bounded derivatives on each packet.
At the initial radial split the barriers are the separate relatively
compact domains in the radial pole bands, and the choices are locally
finite with locally prescribed bounds.
\end{lemma}

\begin{proof}
The nesting of the core closures inside the pole interiors and the
compatible approximation construction of \cref{packet:approximation}
give smooth separating cutoff functions.  A level which is regular on
every packet gives the required domains: there are only countably many
packets, so Sard's theorem permits a common level.  Their far collars
are exact copies of the face domains.  On a fixed packet, compact parts
and the finite face diagram give a positive tubular radius and bounded
geometry for the resulting boundary.  Extend its outward normal in a
smaller tube and multiply by a cutoff of absolute value at most one.
This supplies $Y^\pm$ and a bound for its weighted divergence.  The radial
case uses the relatively compact pole barriers individually.

Fix the barriers and their vector fields first.  By
\cref{packet:approximation}, choose a compatible positive smooth function
$h$ with sufficiently small value and first derivative, and put $M=h^{-1}$.
Require
\[
 M\ge3+|\Div_\omega Y^\pm|\quad\hbox{on }D^\pm,
 \qquad 2|\grad h|\le a_r.
\]
Only finitely many coface requests reach a given packet; at the radial
step the requests are locally finite on compact sets.  The choices
therefore have the stated compatibility.  The derivative inequality implies
$2|\grad M|\le a_rM^2$.  The clipped function
\[
                  z=\max(-M,\min(\mu_{j_0},M))
\]
is finite and locally Lipschitz.  At almost every point its derivative is
that of one active branch, and consequently
\[
                    2|\grad z|\le a_r z^2+\Pi_{j_0}.
\]
This assertion also holds near the original poles, since the clipping
levels are locally bounded and finite truncations of $\mu_{j_0}$ are
locally Lipschitz.  On $D^\pm$ one has $z=\pm M$.

Here is the quantitative smoothing fact being used.  If a locally
Lipschitz function $z$ obeys $|\grad z|\le F$ almost everywhere for a
continuous positive $F$, then, for arbitrary positive continuous error
targets, local convolution in sufficiently small metric charts followed
by a partition of unity gives a smooth $z_1$ with prescribed value error
and $|\grad z_1|\le F$ plus the prescribed error.  Indeed convolution
averages the gradient inequality, the oscillation of $F$ and the metric
on its support tends to zero, and differentiating the partition of unity
produces only the value errors times its fixed derivatives.  A locally
finite cover permits the errors to be assigned summably at each point.
For compatible smoothing, start on deepest faces.  On overlaps in far
collars the clipped functions and all face copies agree.  Blend with the
cutoffs in \cref{packet:approximation}; the same value-error estimate
pays the additional cutoff derivatives.  Preallocate smaller errors on
faces to meet their finitely many coface requests.  Thus the gradient
estimate is always compared with the single continuous function
$(a_rz^2+\Pi_{j_0})/2$, rather than with unrelated bounds for the pieces.

Choose the value error at most
$\min\{1/4,10^{-3}\sqrt{\Pi_{j_0}}\}$ and the gradient error so that
\[
 2|\grad p_0|\le a_rz^2+1.005\Pi_{j_0}.
\]
These are positive continuous error targets.  For $e=p_0-z$,
\[
 .65(z+e)^2-.63z^2\ge-20.475e^2,
\]
so $a_r\le .63$ gives
\[
             2|\grad p_0|\le .65p_0^2+1.01\Pi_{j_0}.
\]
On $D^\pm$ the value error and the earlier choice of $M$ give
$\pm p_0\ge 11/4+|\Div_\omega Y^\pm|$.  Finally set
\[
           c_* =\min\{1/2,.01\sqrt{\pi_{\rm cut}}\}.
\]
The elementary identity
\[
 .67(y+c)^2-.65y^2\ge-21.775c^2\qquad(y,c\in\R)
\]
and $\Pi_{j_0}\ge\pi_{\rm cut}$ prove
\eqref{cut:pressure-gradient} for every $|c|<c_*$; the barrier margin
also persists.  Thus the domains, clipping height, smoothing, and open
shift interval have all been fixed before selecting a minimizing phase.
On a post-radial packet every smooth datum just constructed has bounded
derivatives, by its compact parts and exact products.  At the radial
step the construction gives the corresponding bounds on each compact
window, which is all that is used.
\end{proof}

\begin{definition}[Local absolute minimizer]\label{cut:minimizer-definition}
For smooth $p$ a phase $A$ is a local absolute minimizer if for every
relatively compact open $U\subset E$ and every locally finite-perimeter
set $B$ with $A\mathbin\triangle B\Subset U$,
\begin{equation}\label{cut:local-minimality}
 P_\omega(A;U)-\int_{A\cap U}p\,\dd\mu
 \le P_\omega(B;U)-\int_{B\cap U}p\,\dd\mu.
\end{equation}
The notation $A\mathbin\triangle B\Subset U$ means that the discrepancy
is contained, modulo null sets, in a compact subset of $U$.
\end{definition}

\begin{lemma}[Existence and strict local compactness]\label{cut:bv-compactness}
Every smooth positive weight and smooth locally bounded pressure on a
smooth manifold admit a local absolute minimizer.  Suppose on exhausting
coordinate windows that metrics and positive weights converge smoothly
and pressures converge locally uniformly.  Any sequence of corresponding
local absolute minimizers has a subsequence whose indicators converge in
$L^1_{\rm loc}$ to a local absolute minimizer of the limit data.  Along
that subsequence the weighted perimeter measures converge weakly locally,
as do the unweighted perimeter measures and their associated unoriented
perimeter varifolds.  The assertion includes smooth pointed rescalings
when the rescaled data have the stated convergence.
\end{lemma}

\begin{proof}
We first establish the comparison formulas that will show both
minimality of a limiting phase and absence of perimeter loss. We then
construct exhaustion minimizers and apply these formulas to locally
convergent sequences.

\emph{Comparison formulas.}
Local $L^1$ compactness under local perimeter bounds and lower
semicontinuity are \citet[Chapter~2, Theorem~2.6]{simonGMT}; positive
smooth weights and smooth metrics on a relatively compact chart give
equivalent variation norms. The further BV inputs are the perimeter
lattice inequality and the trace and gluing formulas, both across a
fixed smooth face and across almost every admissible level of a
Lipschitz slicing function. The lattice inequality is
\begin{equation}\label{cut:lattice}
 P_\omega(A\cap B;U)+P_\omega(A\cup B;U)
          \le P_\omega(A;U)+P_\omega(B;U).
\end{equation}
For smooth functions $u,v$, the gradients of $\min(u,v)$
and $\max(u,v)$ are the two original gradients, possibly interchanged,
almost everywhere.  Strictly approximate the two characteristic
functions, apply this identity, and use lower semicontinuity on the
minimum and maximum.  This gives \eqref{cut:lattice}, with its smooth
metric and weight.

We record the trace argument for a fixed smooth comparison face, since
such a face need not be an almost-everywhere selected level.
Work first in a flattened coordinate collar
$Q\times(-a,a)\subset\mathbb R^{m-1}\times\mathbb R$, with the face
$T=Q\times\{0\}$ and the plus side given by $t>0$.
All collars used here are buffered inside a relatively compact ambient
chart.  The functions under consideration are integrable and have
finite total variation on each smaller collar reaching $T$; for the
restrictions of our finite-perimeter phases this follows from local
finite perimeter in the ambient chart.  Local BV regularity only on
sets bounded away from $T$ would not suffice.

Let $w$ be such a BV function on the plus half-collar, and let $D_t w$
denote its distributional derivative in the $t$ direction.
Choose a compact tangential box $K$ and an open box $K'$ with
$K\Subset K'\Subset Q$, and fix $a_0<a$.
Mollify on the open half-collar, using kernels small enough to remain
inside it on each fixed strip.  Fix a countable exhaustion by buffered
tangential boxes and compact height intervals contained in $(0,a)$.
Choose one subsequence whose local $L^1$ errors are summable on every
member of this exhaustion, by making its $k$th error at most $2^{-k}$
on each of the first $k$ members.  Fubini then gives convergence in
$L^1$ on horizontal slices for all heights in one common set of full
measure, simultaneously on the exhaustion boxes.  Restriction gives
the same convergence on $K'$.
For a smooth mollification the fundamental
theorem of calculus, followed by integration over $K$, bounds the
difference of two slices by the integral of the absolute vertical
derivative on the intervening strip.  Convolution bounds this integral
by the variation of $D_t w$ on a slightly enlarged strip; the
tangential enlargement lies in $K'$ when the kernel is small enough.
In the limit the enlarged tangential box lies in $K'$ and the height
interval is contained in $[s,t]$; thus the limiting bound is at most
$|D_t w|(K'\times[s,t])$.
Excluding heights at which
$|D_t w|(K'\times\{t\})>0$, gives
\begin{equation}\label{cut:fixed-face-slices}
 \int_K |w(y,t)-w(y,s)|\,\dd y
 \le |D_t w|\bigl(K'\times(s,t)\bigr),
 \qquad 0<s<t<a_0,
\end{equation}
for all $s,t$ in a common set of full measure.
There are at most countably many excluded atom heights, since the
variation on the buffered collar is finite.

The right-hand side of \eqref{cut:fixed-face-slices} tends to zero
uniformly for good heights $s,t\downarrow0$: the sets
$K'\times(0,\delta)$ decrease to the empty set inside the open
half-collar and have finite variation measure.  The good slices are
therefore Cauchy in $L^1(K)$ and have a limit, denoted
$\operatorname{Tr}_+w$.
Traces obtained on nested tangential boxes agree by uniqueness of the
$L^1$ limit, giving a locally integrable trace on $T$.
Reflection gives the analogous trace $\operatorname{Tr}_-w$ from the
minus side.  For characteristic functions these traces are themselves
characteristic functions almost everywhere, since an $L^1$ limit of
characteristic functions has an almost-everywhere convergent
subsequence.  In particular the phase traces take values in $\{0,1\}$.

The trace also gives the required integration-by-parts formula.
On the plus side, multiply a smooth vector field, compactly supported
in the tangential directions and away from the outer end of the
collar, by a cutoff which is zero for $t\le\varepsilon$ and one for
$t\ge2\varepsilon$.
Apply the distributional derivative identity inside the open
half-collar and let $\varepsilon\downarrow0$.
The cutoff derivative term tends to the pairing of
$\operatorname{Tr}_+w$ with the vector field's $t$ component at $T$:
its derivative has integral one, and the good slices converge in
$L^1$ uniformly as their heights approach zero.
The other terms converge by integrability and finite variation.
This gives the boundary term with outward normal $-\partial_t$ on
the plus side; on the minus side the outward normal is $\partial_t$.
Undoing the smooth coordinate change gives the corresponding
Riemannian formula.  More explicitly, the smooth coordinate volume
density is included in the divergence test; its boundary flux is
exactly $\langle X,\nu\rangle\,\dd\mathcal H_g^{m-1}$.
Thus, if $D_g w$ denotes the intrinsic distributional gradient measure,
with convention
$\int w\Div_g X\,\dd\vol_g=-\int\langle X,\dd D_g w\rangle$
for tests supported in the open side, then
\[
 \int_{U_\pm}w\Div_g X\,\dd\vol_g
 =-\int_{U_\pm}\langle X,\dd D_g w\rangle
   +\int_T\operatorname{Tr}_\pm w\,
                   \langle X,\nu_{\mathrm{out},\pm}\rangle
                         \,\dd\mathcal H_g^{m-1}.
\]
The vector field here has compact support in the collar, so no other
boundary contribution occurs.

Now orient $\nu_T$ from the minus side to the plus side, and glue a BV
function $u$ on the minus side to a BV function $v$ on the plus side.
Adding their integration-by-parts identities shows that the gradient
measure of the glued function is
\[
 (D_g u)\lfloor U_-+(D_g v)\lfloor U_+
 +\bigl(\operatorname{Tr}_+v-\operatorname{Tr}_-u\bigr)
                  \nu_T\,\mathcal H_g^{m-1}\lfloor T.
\]
These three vector measures are concentrated on the disjoint
measurable sets $U_-$, $U_+$, and $T$.  Their variations therefore
add.  Multiplying the variation measures by the smooth positive
weight $\omega$ gives the weighted formula, whose extra interface
cost is precisely
\[
 \int_T\omega\,
       |\operatorname{Tr}_+v-\operatorname{Tr}_-u|
                         \,\dd\mathcal H_g^{m-1}.
\]
We are weighting the variation, not differentiating $\omega$ times
the glued function.  For phases the last integrand is the indicator
of the trace-disagreement set, multiplied by $\omega$.

We will also glue across almost every level of a locally Lipschitz
slicing function $f$, whose level sets need not be smooth.
For this distinct use, the reduced-boundary representation identifies
a phase locally with its integer-multiplicity integration current and
identifies boundary mass with perimeter
\citep[Chapter~3, Section~4, Equation~(4.1) and Theorem~4.3]{simonGMT}.
In a locally oriented chart the phase current has locally finite mass
and locally finite boundary mass.  Lipschitz coarea and the current
slicing and restriction formulas therefore apply
\citep[Chapter~3, Section~2, Equations~(2.8)--(2.9);
Chapter~6, Section~4, Lemmas~4.1 and~4.5 and
Definition~4.4]{simonGMT}.
In particular, writing $J$ for either phase current, for almost every
level $\tau$ one has
\[
 \langle J,f,\tau\rangle
   =\partial\bigl(J\lfloor\{f<\tau\}\bigr)
             -(\partial J)\lfloor\{f<\tau\}.
\]
Apply this identity to both phases, excluding also the null set of
levels charged by their local volume or perimeter measures.
For a top-dimensional phase current the slice has multiplicity one
on the sliced phase and zero elsewhere.  Indeed the slicing
definition restricts its phase indicator to the noncritical part of
the level; coarea shows that critical points and any changes of the
ambient almost-everywhere representative affect almost no level in
positive hypersurface measure.  The two phase slices have the same
orientation induced by $f$ in the locally oriented chart, so their
difference has multiplicity equal to the difference of the two sliced
indicators.
The boundary of their gluing has the two restricted boundary currents
and the difference of their level slices as its interface term.
Its extra weighted mass is bounded by the weighted hypersurface
measure of the set where the sliced phase values differ.
This is the almost-everywhere interface-cost bound used in the
annular comparisons, distance-ball comparisons, and level-set
deletions below.  A countable collection of phases, including a
converging sequence and its limit, can use one common set of
admissible levels by intersecting their full-measure sets.
If the phase indicators converge locally in $L^1$, coarea on a fixed
buffered slicing annulus makes the integrated sliced disagreement
tend to zero.  Pass to a subsequence for which these integrals are
summable.  Fubini then makes the disagreement tend to zero at almost
every level, still within the common admissible set.
This supplies a fixed gluing level with vanishing interface cost in
the compactness argument below.  All orientations in this argument
are local; global orientability is not required.

\emph{Exhaustion minimizers and local compactness.}
Choose a smooth relatively compact exhaustion $U_k\Subset E$.
Minimize over phases vanishing outside $\overline U_k$, taking the total
perimeter, including the interface on $\partial U_k$, and the pressure
integral inside $U_k$.  The empty phase is admissible.  Because the
pressure is bounded on $\overline U_k$, an energy bound gives a perimeter
bound.  The class is closed under $L^1$ convergence, the volume term is
continuous, and perimeter is lower semicontinuous.  The direct method
thus gives a minimizer, locally absolute minimizing in $U_k$.

The same local bound applies to arbitrary local absolute minimizers.
In a coordinate ball $B_t$ compactly contained in their comparison
domain, comparison with the phase cleared inside $B_t$ gives, for almost
every $t$,
\[
 P_\omega(A;B_t)\le\int_{\partial B_t}\omega\,\dd\mathcal H_g^{m-1}
                         +\int_{B_t}|p|\,\dd\mu.
\]
Choose $t$ in a fixed annulus around the smaller window to be bounded.
Coarea bounds the first term for a suitable slice.  Smoothly convergent
metrics and positive weights and locally uniformly bounded pressures
make this a uniform local bound.  Diagonal BV compactness therefore
produces a locally $L^1$ convergent subsequence, both for the exhaustion
minimizers and for any sequence in the statement.

\emph{Minimality and perimeter convergence.}
The gluing formulas now apply to the locally convergent subsequence.
Choose nested relatively compact windows $U\Subset V$
with a Lipschitz foliation of the intervening annulus.  The integral over
slicing parameters of the trace discrepancy between $A_i$ and $A$ tends
to zero by $L^1$ convergence and coarea.  One can therefore choose a
slicing hypersurface between $U$ and $V$ whose discrepancy tends to zero
along a subsequence.  Glue $A$ inside this surface to $A_i$ outside.
Minimality of $A_i$ bounds its interior perimeter from above by that of
$A$, the vanishing trace cost, and a vanishing difference of volume
terms.  The last difference vanishes by the convergence of the data and
indicators.  First use windows whose limiting perimeter gives zero mass
to their boundaries; then shrink the annulus.  Together with lower
semicontinuity this proves weak local convergence of the perimeter
measures.  Replacing the interior copy of $A$ by any compact competitor
$B$ gives the same upper comparison with $B$, since $B=A$ on the gluing
annulus.  Lower semicontinuity on the left proves
\eqref{cut:local-minimality} for $A$.  Smooth metric and weight changes
make only errors tending uniformly to zero on each fixed window, so the
same argument proves the convergent-data assertion.

\emph{Changing metrics and varifold convergence.}
For the varifold assertion, account explicitly for the changing metric.
In a coordinate chart let $D_0\chi_{A_i}$ be the
Euclidean distributional derivative and let $(g_i)$ denote the positive
metric matrices.  The vector measures
\[
      \mathbf v_i=\sqrt{\det g_i}\,g_i^{-1/2}D_0\chi_{A_i}
\]
have Euclidean total variation equal to the $g_i$-perimeter measures.
The matrices and their derivatives converge smoothly, so distributional
convergence and the local mass bounds give $\mathbf v_i\rightharpoonup
\mathbf v$.  The perimeter conclusion gives
$|\mathbf v_i|\rightharpoonup|\mathbf v|$ as well.

This strict convergence also controls the unit polar vectors.  To see
this, on a compact window approximate the limiting polar vector in
$L^2(|\mathbf v|)$ by a continuous bounded vector field $e$.
Expand
\[
 \int|\mathbf v_i/|\mathbf v_i|-e|^2\,\dd|\mathbf v_i|
 =\int(1+|e|^2)\,\dd|\mathbf v_i|-2\int e\cdot\dd\mathbf v_i.
\]
Weak convergence gives its limit, and then the approximation makes that
limit arbitrarily small.  Uniform continuity consequently gives
convergence for every continuous compactly supported integrand of
position and polar vector.  The unoriented $g_i$-orthogonal tangent
hyperplane is a continuous function of this vector and $g_i$.
This proves varifold convergence.  Multiplication or division by the
smooth positive weight on a fixed compact window gives the asserted
weighted and unweighted versions.
\end{proof}

\subsection{Barriers, ordering, and exact product limits}

\begin{lemma}[Barriers hold for every minimizer]\label{cut:barriers}
Every local absolute minimizer for a pressure of
\cref{cut:pressure} contains $D^+$ and excludes $D^-$.  Its boundary
support consequently misses a neighborhood of each current core closure.
After the radial split this neighborhood can be chosen with a positive
uniform radius on each packet.  At the radial split the assertion is
local on the corresponding compact radial blocks.
\end{lemma}

\begin{proof}
We first prove the positive assertion.  Put $D=D^+$, $Y=Y^+$ and
$G=D\setminus A$.  For almost every distance radius $R$ about a fixed
point, $U=B_R$ is a finite-perimeter set, its boundary has zero mass for
the relevant perimeter measures, and BV traces are available.  Write
\[
 V(R)=\mu(G\cap B_R),\qquad
 S_R=\int_{\partial^*B_R}\operatorname{Tr}(\chi_G)\omega\,
                  \dd\mathcal H^{m-1}.
\]
Distance coarea, with $|\grad d|=1$ almost everywhere off the center,
gives local absolute continuity and $V'(R)=S_R$ for almost every $R$.
Weighted Gauss--Green, $Y=\nu_D$ on $\partial D$, and $|Y|\le1$ give
\begin{align*}
 P_\omega(D;U)-P_\omega(A\cap D;U)
 &\le\int_{G\cap U}\Div_\omega Y\,\dd\mu\\
 &\quad-\int_{\partial^*U}
             \operatorname{Tr}(\chi_G)\,Y\cdot\nu_U\,
                         \omega\,\dd\mathcal H^{m-1}.
\end{align*}
By \eqref{cut:lattice} this bounds the interior perimeter increase on
replacing $A$ by $A\cup D$ in $U$.  Glue back to $A$ outside $U$; the
extra perimeter is at most $S_R$.  This is a compact comparison in a
slightly larger relatively compact open set.  By
\eqref{cut:barrier-pressure} and minimality,
\begin{equation}\label{cut:barrier-volume}
                         V(R)\le2S_R=2V'(R).
\end{equation}
The weighted ball volume is at most polynomial by
\cref{packet:geometry}, since $\omega$ is bounded on the packet.
If $V(R_0)>0$, \eqref{cut:barrier-volume} would imply
$V(R)\ge V(R_0)e^{(R-R_0)/2}$, a contradiction.  Thus $G$ is null.
For a relatively compact radial barrier, simply add all of $D$; no
slicing terms occur, and the same comparison gives $\mu(G)=0$ directly.
The complement of $A$ minimizes pressure $-p$, proving the negative
assertion.  The domains contain the core closures in their interiors;
the packetwise positive separation follows from exact products and the
compact-part property.  A phase equal to one or zero on an open set has
no perimeter support there.
\end{proof}

\begin{proposition}[Generic uniqueness on all required products]
\label{cut:generic-uniqueness}
After the radial split one common shift $c$, with $|c|<c_*$, can be chosen
such that there is exactly one local absolute minimizer, modulo null
sets, on every packet $E$ and every product $E\times\R^d$, $d\ge0$, for
the pulled-back weight and pressure $p_0+c$.  The unique product phase
is exactly $A_E\times\R^d$ modulo null sets, where $A_E$ is the unique
phase on $E$.
\end{proposition}

\begin{proof}
On any one such space, consider \emph{arbitrary} minimizers $A$ for
$p_0+c_1$ and $B$ for $p_0+c_2$, with $c_1<c_2$.  Put
$G=A\setminus B$, $\delta=c_2-c_1$, and use the same distance-ball
slicing as in the preceding proof.  Replace $A$ by $A\cap B$ inside the
ball and $B$ by $A\cup B$ there, gluing each to its own exterior phase.
The two gluing costs total at most $2S_R$.  Their interior perimeter
increments have nonpositive sum by \eqref{cut:lattice}; their pressure
increments sum to precisely $-\delta\mu(G\cap B_R)$.  Adding the two
minimality inequalities gives
\begin{equation}\label{cut:ordering-volume}
                 \delta V(R)\le2V'(R).
\end{equation}
Polynomial weighted growth, also on $E\times\R^d$, forces $V=0$.
Thus \emph{every} minimizer at the smaller shift is contained in
\emph{every} minimizer at the larger shift.

Choose countably many nonnegative compactly supported continuous test
functions $\varphi_a$ which separate locally integrable functions.
For each $a$ and shift $c$, let
\[
 I_a(c)=\left(\inf_{A\text{ minimizes at }c}
                 \int\varphi_a\chi_A\,\dd\mu,
              \ \sup_{A\text{ minimizes at }c}
                 \int\varphi_a\chi_A\,\dd\mu\right).
\]
Existence is supplied by \cref{cut:bv-compactness}.  The preceding
ordering says that the nonempty open intervals $I_a(c)$ for distinct
$c$ are disjoint.  Each contains a rational number, so only countably
many can be nonempty.  If minimizers differ, some $\varphi_a$ distinguishes
them and $I_a(c)$ is nonempty.  There are therefore at most countably
many exceptional shifts on this space.  The packet collection and the
set of integers $d$ are countable, so their exceptional shifts have a
countable union.  Choose $c$ in the remaining part of $(-c_*,c_*)$.

A Euclidean translation of the unique minimizer on $E\times\R^d$ is
again a minimizer.  Uniqueness implies translation invariance of its
indicator.  It suffices to impose the countable dense group of rational
translations: continuity of translations in $L^1_{\rm loc}$ then gives
all translations, so its Euclidean distributional derivatives vanish.
The indicator is therefore $\chi_F(x)$ for a base phase $F$.
BV slicing shows that $F$ has locally finite perimeter.  If a compact
base competitor improved its energy by $\eta>0$, use that competitor
over $(-L,L)^d$ and the phase $F$ outside, retaining the original phase
outside the compact base discrepancy.  The energy gain is
$\eta(2L)^d$.  The interfaces at the Euclidean box boundary cost at most
$O(L^{d-1})$ times the weighted volume of that compact discrepancy.
For large $L$ this contradicts minimality.  Thus $F$ minimizes on $E$,
and base uniqueness identifies it with $A_E$.  The case $d=0$ needs
no product argument.
\end{proof}

At the radial split set $c=0$ and choose any minimizer supplied by
\cref{cut:bv-compactness}.  Its compact barriers confine its boundary to
the locally finite compact transition shells.  No uniqueness or
polynomial-growth assertion on the initial manifold is used.  We call
the selected minimizing boundary a \emph{true cut}; it will later be
replaced by a compatible regular wall near the required region.

\subsection{The precise local regularity input}

\begin{theorem}[Local GMT properties used here]\label{cut:gmt}
Let $A$ satisfy \eqref{cut:local-minimality} on a smooth $m$-manifold,
$m\ge2$, with smooth positive $\omega$ and smooth $p$.  Put
$S=\supp|D\chi_A|$.  Then $S=S_{\rm reg}\cup\Sigma$, where $S_{\rm reg}$
is a smooth embedded hypersurface, $\Sigma$ is relatively closed, and
\begin{equation}\label{cut:singular-dimension}
                         \dim_{\mathcal H}\Sigma\le m-8.
\end{equation}
A negative dimension bound means that $\Sigma$ is empty.  The outward
measure-theoretic normal extends smoothly on $S_{\rm reg}$.  With
$B(u,v)=\langle\nabla_u\nu,v\rangle$, its base mean curvature satisfies
\begin{equation}\label{cut:base-mean-curvature}
                    H_g=p-\partial_\nu\log\omega.
\end{equation}
The unweighted multiplicity-one perimeter varifold has locally bounded
generalized mean curvature.  Its density exists at every support point;
monotonicity gives local uniform area upper bounds and positive area
lower bounds at support points.  Under the convergences of
\cref{cut:bv-compactness}, densities are upper semicontinuous at
converging points, provided the local generalized-mean-curvature bounds
are uniform.  The supports converge locally, and on every compact
subset of the regular limiting hypersurface convergence is smooth and
single-sheeted when the data converge smoothly.  On compact subsets of
the interiors of the limiting phases there is eventually no boundary.

Every tangent at a point of $S$ is a multiplicity-one locally
perimeter-minimizing boundary cone in the Euclidean tangent space.
A flat tangent implies regularity.  A tangent at a singular point is
therefore nonflat, and its translation-invariance spine has dimension
at most $m-8$.
\end{theorem}

\begin{proof}[Proof by first variation, strict blowups, and regularity]
Put $d=m-1$.  The classical inputs are the bounded-mean-curvature
monotonicity theorem and Allard's graphical theorem
\citep[Chapter~4, Theorem~3.17; Chapter~5, hypotheses~5.1 and
Theorem~5.2]{simonGMT}, and the singular-dimension theorem for
\emph{exactly} minimizing boundary currents
\citep[Chapter~7, Theorem~5.8]{simonGMT}.  The minimizing-boundary dimension bound follows from the work of
Simons and Federer \citep{simons1968,federer1970}. The last theorem will be applied
only to constant-metric blowup limits.  We verify the transfer to the
present smooth weighted problem, including the singular-set compactness
needed for the Hausdorff-dimension conclusion.

\emph{First variation and local area bounds.}
First variation is available before any regularity conclusion.  The
flow of a compactly supported smooth field $Z$ is an admissible
comparison in \eqref{cut:local-minimality}, so
\[
 \int_{\partial^*A}\omega
   \bigl(\Div_S Z+\langle\grad\log\omega,Z\rangle\bigr)
       \,\dd\mathcal H_g^d
 =\int_{\partial^*A}p\omega\langle Z,\nu\rangle
       \,\dd\mathcal H_g^d.
\]
Here tangent planes and $\nu$ are the measure-theoretic ones.  Substituting
$Z=W/\omega$ cancels the tangential weight derivatives and gives
\[
 \int_{\partial^*A}\Div_S W\,\dd\mathcal H_g^d
  =\int_{\partial^*A}
       (p-\partial_\nu\log\omega)\langle W,\nu\rangle
          \,\dd\mathcal H_g^d.
\]
Thus the unweighted perimeter varifold has intrinsic generalized
mean-curvature vector
$-(p-\partial_\nu\log\omega)\nu$, locally bounded by
$|p|+|\grad\log\omega|$.  Its multiplicity is one almost everywhere.

On a relatively compact coordinate patch, first pass to a smaller
patch and extend its metric to a closed manifold.  Explicitly, place
its coordinate closure inside a coordinate ball of a flat torus and
blend the given positive definite metric with the flat metric outside
the smaller patch; the convex combination stays positive definite.
The resulting metric is smooth and the ambient manifold is compact,
so the smooth isometric embedding theorem applies
\citep[Theorem~2]{nash1956}.  The preceding intrinsic first variation then gives
Euclidean generalized mean curvature with the additional bounded term
$\tr_{TS}\mathrm{II}$, where $\mathrm{II}$ is the second fundamental
form of this fixed embedding.  Indeed decompose a Euclidean test field
into its ambient-manifold tangent and normal parts and differentiate
that decomposition along $TS$.  Its normal contribution is exactly the
trace of $\mathrm{II}$.  Consequently Euclidean monotonicity applies
locally.  Intrinsic and embedded density limits agree, since the
embedding is isometric to first order.  Monotonicity gives density
existence and upper semicontinuity, including at nonreduced support
points; their densities are at least one, by approximation by reduced
points.  Together with the clearing-ball upper bound from
\cref{cut:bv-compactness}, it gives uniform upper and positive lower
area bounds on compact patches.

\emph{A common ambient model for varying metrics.}
We record carefully the uniform version when the metric changes.
Suppose $g_i\to g$ smoothly on a fixed patch, and all other data have
the local bounds in the theorem.  Any claimed failure of regular-point
persistence or support convergence has a subsequence on which the
failure persists.  Pass further so that, on a slightly larger compact
patch,
\[
                  \|g_i-g\|_{C^i}\le 2^{-i^2}.
\]
Set $t_i=2^{-i}$ and take disjoint smooth bump functions $\beta_i(t)$,
equal to one on a neighborhood of $t_i$, supported in intervals of
width comparable to $t_i$, and satisfying
$|\beta_i^{(b)}|\le C_b t_i^{-b}$.  After omitting finitely many terms,
\[
 g_t=g+\sum_i\beta_i(t)(g_i-g),\qquad G=\dd t^2+g_t
\]
is a smooth positive metric on a product neighborhood, with $g_0=g$
and $g_{t_i}=g_i$.  For each pair of fixed orders $a,b$, the summands
after $a$ spatial and $b$ temporal derivatives are bounded by
$C_b2^{ib-i^2}$ for all sufficiently large $i$.  These bounds, divided
by any fixed power of $t_i$, still tend to zero.  Thus all mixed
derivatives extend at $t=0$.  The plateaus also give
$\partial_tg_t=0$ at $t_i$, so the slices carrying $g_i$ are totally
geodesic in $G$.

Extend this \emph{one} smooth metric $G$ across the edge of a smaller
product patch to a closed manifold and isometrically embed it once.
The images of the slices $t=t_i$ converge smoothly to the slice
$t=0$, and the second fundamental form of this single embedding is
bounded on their common compact neighborhood.  The preceding
first-variation calculation therefore gives one Euclidean
mean-curvature bound for all the embedded perimeter varifolds.
Strict perimeter convergence from \cref{cut:bv-compactness} becomes
varifold convergence in this same Euclidean space.  This construction
proves the needed sequential uniformity; it does not assume that
independently chosen isometric embeddings have uniform bounds.

\emph{Exact minimizing tangent cones.}
Next consider a blowup of one phase at $x\in S$.  In smooth coordinates,
at scale $\rho$, the normalized data are
\[
 A_{x,\rho}=\rho^{-1}(A-x),\quad
 g_{x,\rho}(y)=g(x+\rho y),\quad
 \omega_{x,\rho}(y)=\omega(x+\rho y),\quad
 p_{x,\rho}(y)=\rho p(x+\rho y).
\]
The metric here has been divided by the square of the physical scale:
perimeter scales by $\rho^d$, volume by $\rho^{d+1}$, and this accounts
for the factor $\rho$ in the pressure.  On every fixed window the data
converge smoothly to $g(x),\omega(x),0$.  The strict compactness lemma
therefore identifies every subsequential tangent varifold with the
unit-multiplicity perimeter varifold of a phase minimizing perimeter
in the constant metric $g(x)$.  A constant weight can be divided out.
After one linear isometry this is an ordinary Euclidean
perimeter-minimizing phase.  It is also a minimizing boundary current.
Indeed, orient this Euclidean chart.  A compactly supported integral
current filling a boundary competitor has top dimension and is given
by an integer-valued BV function $q$.  The competitor is the boundary
of $\chi_A+q$.  Clipping that integer-valued function to $[0,1]$ leaves
it unchanged outside the compact discrepancy and gives a characteristic
function without increasing variation.  Any current competitor with
smaller mass would therefore give a smaller-perimeter set competitor.
This verifies the precise minimizing-boundary hypothesis of
\citet[Chapter~7, Theorem~5.8]{simonGMT}.  All orientations in this
argument are local.

The support lower bound makes this limit nonzero at the origin.
Monotonicity based at $x$ makes its mass ratio constant on all radii;
the equality case, or the stationary radial identity in
\citet[Chapter~4, Equation~(3.6)]{simonGMT}, makes it conical.
Equivalently, after a fixed local isometric embedding, dilations of
the embedded ambient manifold converge smoothly to its tangent
space, its second fundamental form is multiplied by $\rho$, and
the same blowups give this cone.  Its multiplicity is one as a
perimeter varifold, although its density at the vertex need not
be one.

\emph{Regularity and persistence under convergence.}
We now apply Allard's graphical regularity theorem \citep{allard1972}.  In its Euclidean form used here,
fix an exponent $q>d$.  For a rectifiable $d$-varifold of multiplicity
at least one, on a ball centered at a support point and contained in
the first-variation domain, sufficiently small excess of the normalized
mass above one and sufficiently small
\[
       \rho^{\,1-d/q}
       \left(\int_{B_\rho}|\mathbf H|^q\,\dd\mu_V\right)^{1/q}
\]
give a single $C^{1,1-d/q}$ graph on a smaller ball
\citep[Chapter~5, Theorem~5.2]{simonGMT}.  A unit flat tangent supplies
the small mass excess.  The bounded mean curvature supplies the other
small quantity, since it is $O(\rho)$ under the local area bound.
Thus a flat tangent makes the original point regular as a
$C^{1,\alpha}$ graph.  The graph satisfies the weak prescribed-curvature
equation
$H_g=p-\partial_\nu\log\omega$; its coefficients and right-hand side
are smooth functions of position, graph height, and graph gradient.
The bounded slope gives uniform ellipticity.  We spell out the passage
from a weak graph to a smooth one.  In a relatively compact chart write
the equation as
\[
 \partial_i A^i(x,u,Du)=B(x,u,Du).
\]
Here $A$ is the gradient, in the slope variables, of the smooth
positive weighted graph-area integrand.  Thus $A_p$ is symmetric and
uniformly positive on the bounded slope range; $B$ is smooth in its
arguments.  For a compactly supported cutoff $\eta$, test translated
differences with $\eta^2\delta_hu$.  The difference of $A$ has the form
$a_hD\delta_hu+c_h\delta_hu+b_h$, where $a_h$ is uniformly elliptic
and all three coefficients are uniformly bounded.  Since $u$ is
Lipschitz, $\delta_hu$ is bounded.  On the right move the difference
quotient from $B$ to the test function, using
$\|\delta_{-h}\phi\|_2\le\|D\phi\|_2$.
Ellipticity and Cauchy's inequality then give
\[
 \int\eta^2|D\delta_hu|^2\le C_\eta
\]
independently of small $h$.  Hence $u\in H^2_{\mathrm{loc}}$.
For $v=\partial_ku$, weak differentiation now gives
\[
 \partial_i(a^{ij}\partial_jv)=\partial_iF^i,\qquad
 a^{ij}=A^i_{p_j},\qquad
 F^i=\delta_{ik}B-A^i_{x_k}-A^i_uv,
\]
with coefficients evaluated at $(x,u,Du)$.  Both $a$ and $F$ are
$C^{0,\alpha}$ because $u$ is $C^{1,\alpha}$.

To justify the divergence Schauder estimate for this weak $H^1$
solution, mollify $a,F$ on a slightly larger ball and solve the
coercive linear Dirichlet problems on the ball with the same trace as
$v$.  Their solutions $v_\epsilon$ satisfy
\[
 \lambda\|D(v_\epsilon-v)\|_2
 \le \|F_\epsilon-F\|_2
       +\|a_\epsilon-a\|_\infty\|Dv\|_2,
\]
where $\lambda>0$ is a common ellipticity constant.  Thus they converge
in $H^1$.  Smooth-coefficient Sobolev regularity makes each approximant
smooth in the interior; the local $L^2$--$L^\infty$ estimate and then
the a priori divergence Schauder estimate give uniform interior
$C^{1,\alpha}$ bounds
\citep[Lecture~6, Theorem~1 and Corollary~1; Lecture~18, Theorem~2;
Lecture~12, Theorem~1]{simonPDE}.
These bounds use only the common ellipticity and H\"older bounds,
not higher derivatives of the mollified coefficients.  Passing to
the limit gives $v\in C^{1,\alpha}_{\mathrm{loc}}$, hence
$u\in C^{2,\alpha}_{\mathrm{loc}}$.  Repeating the difference-quotient
and weak Schauder argument for differentiated equations makes the
graph smooth.  The estimates are uniform for families with common
$C^{1,\alpha}$ graph bounds and uniformly smooth geometric data.

Likewise, near a regular point of a limiting phase choose a sufficiently
small ball with mass ratio as close to one as needed, with strict slack
in the Allard threshold, and with small scaled mean-curvature bound.
Choose its radius also so that its sphere has zero limiting perimeter
measure; almost every sufficiently small radius has this property.
For smoothly convergent data, use the single ambient embedding
construction above.  Weak perimeter-measure convergence now passes the
mass bound on this ball to the approximating phases.  To center at a
nearby approximating support point, decrease the radius by a fixed
arbitrarily small factor: the recentered ball lies in the chosen ball
for all large indices, and the strict slack pays the change in the
normalizing power of the radius.  Allard therefore gives one graph in
each smaller ball for all sufficiently large indices.  The graph equation and smooth data convergence then give
smooth convergence on smaller compact patches.  This also proves
\[
 x_i\in\Sigma_i,\quad x_i\longrightarrow x
                    \quad\Longrightarrow\quad x\in\Sigma
\]
for interior limits: a regular limiting point has just been excluded,
and a point outside the limiting support is excluded by the uniform
support lower bound and perimeter convergence.  The same lower bound
and weak convergence prove local support convergence; they also exclude
stray support in compact subsets of phase interiors.  The regular set
is open in the support, so its complement $\Sigma$ is relatively closed.

\emph{Transfer of the singular-dimension bound.}
For the singular-dimension transfer we use the content argument in
\citet[Chapter~1, Theorem~3.26(2), and Appendix~A, proof of Theorem~A.4,
step~(1)]{simonGMT}.
Its hypotheses in the present application have all been established:
on every compact coordinate patch and at every center, each vanishing
scale sequence has a strict tangent subsequence; these tangents are
homogeneous minimizing boundaries; singular sets commute exactly with
translations and dilations; and, on compact windows of every such
convergent sequence, singular sets eventually lie in every open
neighborhood of the limiting singular set.  For the last assertion,
a contrary sequence of points has a convergent subsequence, contradicting
regular-point persistence or support convergence.  The data on the
rescaled windows have uniform local derivative bounds, with derivatives
of order $j$ acquiring factors $\rho^j$ and pressure derivatives acquiring
$\rho^{j+1}$.  Thus the compactness and regularity argument applies to
all these sequences.  We use only the cited content-transfer step, and
spell it out next; no identification with the scalar-function family
of Appendix~A is needed.

Suppose $m\ge8$ and, for some
$s>m-8$, the singular set in an interior window had positive
$\mathcal H^s$ measure.  Let $\mathcal H^s_\infty$ denote Hausdorff
content, defined by covering with balls and summing their radii to
the power $s$; normalization constants are immaterial.
The Hausdorff-content density selection in the proof of Theorem~A.4
provides an interior $x\in\Sigma$, scales $\rho_i\downarrow0$, and
$c>0$ such that
\[
       \rho_i^{-s}\mathcal H^s_\infty
               (\Sigma\cap\overline B_{\rho_i}(x))\ge c.
\]
Pass to a strict tangent limit $C$ along these scales.  The exact
minimizing-boundary theorem gives
$\mathcal H^s(\Sigma_C)=0$, because its hypersurface dimension is
$d=m-1$ and its bound is $d-7=m-8$.
Cover $\Sigma_C\cap\overline B_{3/2}$ by finitely many open balls
whose sum of $s$th powers of radii is less than $c/2$; the cover can
first be chosen countable with this bound and then reduced to a finite
subcover by compactness.  Singular-set persistence puts every
$\rho_i^{-1}(\Sigma-x)\cap\overline B_1$ in that same union for large
$i$: otherwise a convergent sequence in its compact complement would
give a singular limiting point outside the cover.  Scaling of
Hausdorff content now contradicts the preceding lower bound.
This proves $\mathcal H^s(\Sigma)=0$ for every $s>m-8$.
It is the density selection and this finite-cover passage, in addition
to tangent information, that transfer the dimension bound.

If $m\le7$, the exact minimizing-boundary theorem says that every tangent
cone is regular at its vertex and hence is a hyperplane.  Flat-tangent
regularity then makes $\Sigma$ empty.  Finally a nonflat minimizing
cone is singular at its vertex, and translation invariance carries
that vertex to each point of its spine.  The spine is therefore
contained in its singular set, giving the asserted spine bound.
The preceding monotonicity argument, applied in the common embeddings
for convergent data, also gives upper semicontinuity of densities at
moving points by first passing to the limit at a fixed radius and then
letting the radius decrease to zero.  The smooth graph's coorientation
comes from the phase, and its first-variation equation is
\eqref{cut:base-mean-curvature}.  All constructions are local in the
base and require no ambient orientability or torus-dimensional
regularity theorem.
\end{proof}


\begin{corollary}[True cuts converge at packet ends]
\label{cut:end-convergence}
For the common generic shift, let $A_E$ be the true phase on each packet
and let $A_B$ be that on a face.  On translated windows escaping in a
product collar over a compact part of $B$, the phases $A_E$ converge to
$A_B\times\R^d$ in $L^1_{\rm loc}$, their perimeter measures and
varifolds converge, and their regular boundaries converge smoothly on
compact sets off the limiting singular set.  Their boundary supports
converge locally as well.  The same conclusions hold after passage to
any deeper product furnished by \cref{packet:escape}.
\end{corollary}

\begin{proof}
The metric, weights, and shifted pressure in the translated expanding
windows are exactly the data pulled from $B$.  Local compactness
produces a subsequential limiting phase on $B\times\R^d$, and it is
an absolute minimizer by \cref{cut:bv-compactness}.  The uniqueness and
factorization of \cref{cut:generic-uniqueness} identify it with
$A_B\times\R^d$.  Every subsequence has this same possible limit, which
proves convergence of the full sequence.  Apply \cref{cut:gmt} for the
support and smooth assertions.  The deeper-face statement follows by
using the commuting product coordinates in \cref{packet:escape}.
\end{proof}

\subsection{Stability on a wall and the retained-side estimate}

We now choose the representative of each phase which is open away from
$S$ and agrees near $S_{\rm reg}$ with the corresponding smooth closed
side when a trace is required.  Write $S=S_{\rm reg}$ in integrals below.
Every test has compact support away from $\Sigma$; joined-side tests
also avoid the cylinder over $\Sigma$.  The old torus variables remain
invariant, and $\dd\mu_S=\omega\,\dd\mathcal H_g^{m-1}$ is the wall
measure.  The intrinsic scalar curvature $R_S$ and the second
fundamental form $B$ refer to $S\times\T^q$ with its induced old
stabilized metric.  Its torus principal curvature in direction $i$ is
$\partial_\nu\log f_i$, so \eqref{cut:base-mean-curvature} gives
\[
          H=\tr B=H_g+\partial_\nu\log\omega=p
\]
for the outward normal of the plus phase.  The invariant wall Laplacian
is $\omega^{-1}\Div_S(\omega\grad_S\,\cdot\,)$, so its integration by
parts uses precisely $\dd\mu_S$.

\begin{lemma}[The wall stability form]\label{cut:wall-stability}
For every compactly supported smooth invariant test $v$ on $S$,
\begin{equation}\label{cut:wall-form}
 \int_S\bigl(2|\grad v|^2+R_Sv^2\bigr)\,\dd\mu_S
 \ge\int_S\bigl(R+p^2+|B|^2+2\partial_\nu p\bigr)v^2\,\dd\mu_S.
\end{equation}
In particular its right side is at least
$\int_S(R-1.02\Pi_{j_0})v^2\,\dd\mu_S$.
\end{lemma}

\begin{proof}
Weighted perimeter is the area of the old torus product for invariant
phases.  The invariant normal variation with speed $v$ is compactly
supported in the smooth regular part and comes from an admissible base
variation.  Thus its energy has a minimum at parameter zero.
For completeness the variation identities hold in every dimension:
differentiating the induced metric gives $\dot g_{ab}=2vB_{ab}$,
differentiating the unit normal gives $\dot\nu=-\grad_Sv$, and
commuting the two ambient covariant derivatives in the derivative of
$B_{ab}$ gives, after tracing and differentiating the inverse metric,
\[
 \dot H=-\lap_S v-(\Ric_{\widehat g}(\nu,\nu)+|B|^2)v.
\]
The area derivative is $Hv\,\dd\mu_S$, while the pressure-volume
variation is $pv\,\dd\mu_S$.  The first energy derivative is therefore
$\int_S(H-p)v\,\dd\mu_S$.  At $H=p$ the derivatives of the area
measure and of the chosen speed multiply zero, and
$\dot p=v\partial_\nu p$.  Hence the nonnegative second derivative is
\[
 \int_S\left[-v\lap_Sv-
 (\Ric_{\widehat g}(\nu,\nu)+|B|^2+\partial_\nu p)v^2\right]
                         \,\dd\mu_S.
\]
Weighted integration by parts gives
\[
 \int_S|\grad v|^2\,\dd\mu_S
 \ge\int_S(\Ric_{\widehat g}(\nu,\nu)+|B|^2+
                           \partial_\nu p)v^2\,\dd\mu_S.
\]
Compare the first variation in \citet[Lemma~13]{chodoshLi} and the
second-variation calculation in the proof of its Lemma~14. The displayed
calculation proves the identity in the present stabilized setting; the
stability and Gauss-equation mechanism is that of
\citet[Section~1]{schoenYau1979Structure}.
The derivation here uses neither Chodosh and Li's dimension restriction
for smooth existence nor a regularity assertion for the torus product.
Only the already smooth base patch and its invariant variations enter.
The Gauss identity
\[
 R_S=R-2\Ric_{\widehat g}(\nu,\nu)+H^2-|B|^2
\]
gives \eqref{cut:wall-form}.  By \eqref{cut:pressure-gradient},
\[
 p^2+2\partial_\nu p
       \ge p^2-2|\grad p|
       \ge .33p^2-1.02\Pi_{j_0},
\]
which proves the final statement.
\end{proof}

For either phase side $\Omega$, put $p_s=p$ on the plus side and
$p_s=-p$ on the minus side.  The side is itself an absolute minimizing
phase for the signed pressure $p_s$, and its outward stabilized mean
curvature is $H_{\rm out}=p_s$.

\begin{lemma}[Deletion and the seam estimate]\label{cut:side-estimate}
Every nonnegative compactly supported smooth function $f$ on the
closed regular side, with support avoiding the singular set, satisfies
\begin{equation}\label{cut:deletion}
       \int_S f\,\dd\mu_S
       \le\int_\Omega\bigl(|\grad f|+p_sf\bigr)\,\dd\mu.
\end{equation}
For every smooth real test $v$ of the same type,
\begin{equation}\label{cut:side-form}
 3.05\int_\Omega|\grad v|^2\,\dd\mu
       +2\int_S H_{\rm out}v^2\,\dd\mu_S
       \ge-1.02\int_\Omega\Pi_{j_0}v^2\,\dd\mu.
\end{equation}
Both statements extend to compactly supported Lipschitz tests by
approximation.
\end{lemma}

\begin{proof}
For almost every $t>0$, delete $\Omega\cap\{f>t\}$ from the minimizing
phase.  The removed boundary area is the area on $S\cap\{f>t\}$;
the newly introduced boundary is the level surface of $f$ in the
interior of $\Omega$.  The pressure-volume term increases by
$\int_{\Omega\cap\{f>t\}}p_s\,\dd\mu$.  Minimality gives the
corresponding perimeter inequality.  Integrating in $t$, the layer-cake
formula and weighted coarea give \eqref{cut:deletion}.  All comparisons
are compact and occur on regular patches; the perimeter slicing formula
justifies the same argument at exceptional critical levels by discarding
a null set of levels.

Apply \eqref{cut:deletion} to $f=2(-p_s)_+v^2$, initially by Lipschitz
approximation.  Its right side is bounded above by
\[
 \int_{\Omega\cap\{p_s<0\}}
       \bigl[(2|\grad p_s|-2p_s^2)v^2
                         +4|p_s|\,|v|\,|\grad v|\bigr]\,\dd\mu.
\]
Young's inequality gives
\[
 4|p_s|\,|v|\,|\grad v|
        \le3.05|\grad v|^2+\frac4{3.05}p_s^2v^2.
\]
The pressure estimate holds unchanged for $p_s$, and
\begin{equation}\label{cut:young-margin}
                2-\frac4{3.05}-.67=\frac{113}{6100}>0.
\end{equation}
Discard the resulting nonpositive multiple of $p_s^2v^2$.  Thus
\[
 2\int_S(-p_s)_+v^2\,\dd\mu_S
 \le3.05\int_\Omega|\grad v|^2\,\dd\mu
                  +1.02\int_\Omega\Pi_{j_0}v^2\,\dd\mu.
\]
Since $2p_s\ge-2(-p_s)_+$ on the boundary and
$H_{\rm out}=p_s$, this is \eqref{cut:side-form}.
\end{proof}

\subsection{Raw joined forms and the numerical reserve}

Attach $S\times[0,\infty)_t$ to the regular boundary of either side,
freezing all induced metric and old warp data on the cylinder.  This
is a raw joined space: it is smooth on its two pieces, and the seam
retains its scalar corner contribution.  An admissible test is a pair
$v=(v_b,v_c)$, smooth separately up to regular seam patches, compactly
supported away from singularities and their cylinders, satisfying
$v_b|_S=v_c|_{t=0}$.  Compact Lipschitz versions are allowed by local
approximation in regular seam charts.  For the wall alone the tests are
compact smooth functions on $S$ away from $\Sigma$.

Set $C=3.99$.  On a wall define
\[
 Q_S(v)=\int_S(C|\grad v|^2+R_Sv^2)\,\dd\mu_S.
\]
On a joined side define
\begin{align}\label{cut:raw-form}
 Q_\Omega(v)
  ={}&\int_\Omega(C|\grad v_b|^2+Rv_b^2)\,\dd\mu
       +2\int_S H_{\rm out}v_b^2\,\dd\mu_S\notag\\
    &+\int_0^\infty\!\int_S
        \bigl(C(|\grad_Sv_c|^2+|\partial_tv_c|^2)
                         +R_Sv_c^2\bigr)\,\dd\mu_S\,\dd t.
\end{align}
The sign in the seam term can also be checked directly: in outward
arclength $s$ the mean curvature changes from $H_{\rm out}$ on the
bulk side to zero on the frozen cylinder, and the scalar expression
contains $-2\partial_sH$.  Its distributional contribution is
$+2H_{\rm out}$ at the seam.  The smooth realization and its form
convergence will be established in \cref{sec:regularization}.

\begin{proposition}[Raw scalar domination]\label{cut:raw-domination}
Let $V=R-1.02\Pi_{j_0}$, restricted to the wall and pulled back to the
cylinder when necessary.  Then
\begin{align}\label{cut:wall-domination}
 Q_S(v)&\ge\int_S Vv^2\,\dd\mu_S
                    +1.99\int_S|\grad v|^2\,\dd\mu_S,\\
 Q_\Omega(v)&\ge\int_\Omega Vv_b^2\,\dd\mu
       +\int_0^\infty\!\int_S Vv_c^2\,\dd\mu_S\,\dd t\notag\\
 &\quad+.94\int_\Omega|\grad v_b|^2\,\dd\mu
       +1.99\int_0^\infty\!\int_S|\grad v_c|^2
                                      \,\dd\mu_S\,\dd t.
                                      \label{cut:side-domination}
\end{align}
Here $|\grad v_c|^2=|\grad_Sv_c|^2+|\partial_tv_c|^2$.
\end{proposition}

\begin{proof}
The wall assertion follows by adding $(C-2)|\grad v|^2$ to
\cref{cut:wall-stability}.  For the bulk use
\cref{cut:side-estimate}; the remaining gradient coefficient is
$C-3.05=.94$.  Apply the wall assertion to every cylinder slice and
integrate in $t$.  The horizontal coefficient is $C-2=1.99$, while the
vertical coefficient is still $C=3.99$, so lowering the latter to $1.99$
gives \eqref{cut:side-domination}.  The seam trace in the bulk estimate
is precisely the matching trace in \eqref{cut:raw-form}.
\end{proof}

\begin{definition}[Reserved form and test completion]\label{cut:reserve}
For the wall or a joined side define, respectively,
\begin{align*}
 \mathcal N_S(v)&=1.7\int_S|\grad v|^2\,\dd\mu_S,\\
 \mathcal N_\Omega(v)
   &=.8\int_\Omega|\grad v_b|^2\,\dd\mu
       +1.7\int_0^\infty\!\int_S|\grad v_c|^2
                                            \,\dd\mu_S\,\dd t.
\end{align*}
We write $\mathcal N$ when the tested space is understood.  The squared norm
$\|v\|_2^2$ is the sum of the corresponding piecewise weighted integrals.
With
\[
                  e_0=(\lambda_r-\lambda_{r+1})/100>0,
 \qquad \mathfrak h(v)=\mathcal N(v)+e_0\|v\|_2^2,
\]
the test completion means completion of the admissible compact tests
just specified in the norm $\mathfrak h^{1/2}$.  In particular its seam
condition is matching traces, with no zero trace imposed there.
\end{definition}

The coefficients have a reserve even after the later shielding charge
$.9\mathfrak h$.  Specifically,
\begin{equation}\label{cut:remaining-gradient}
                       .94-.9(.8)=.22,
             \qquad 1.99-.9(1.7)=.46.
\end{equation}
The potential term in \cref{cut:raw-domination}, combined with the
parent inequality, satisfies
\begin{equation}\label{cut:localized-potential}
 R-1.02\Pi_{j_0}
 \ge\lambda_r+b_r\sum_{j\ne j_0}\Pi_j
                         +(b_r-1.02)\Pi_{j_0}
                       \qquad\hbox{on }D_E.
\end{equation}
Here $b_r\ge1.07>1.02$.  The raw form inequalities themselves hold for
all admissible compact tests; the replacement of their potential by
the right side of \eqref{cut:localized-potential} is permitted only
where the test is supported over the old required region.
The count inclusions of \cref{packet:counts} supply exactly this
localization for the later required regions.  These are the forms and
constants used in shielding and regularization.  The minimizing cut,
its pressure, and the old weight remain fixed while the other profiles
and their cores are changed.
